% =============================================================================
% 8.2 CONFIGURACIÓN DE ENTORNOS DE EJECUCIÓN
% =============================================================================
\section{Configuración de Entornos de Ejecución}
\label{sec:entornos}

EthosHub opera en dos entornos con topologías distintas: en \textbf{desarrollo}, frontend
y backend corren como procesos separados comunicados por un \textit{proxy}; en
\textbf{producción}, ambos se fusionan en un único \textit{JAR} monolítico servido desde un
contenedor. La figura~\ref{fig:entornos} contrasta ambas topologías.

\vspace{0.3cm}
\begin{figure}[h!]
\centering
\begin{tikzpicture}[node distance=0.6cm and 1.1cm,
    box/.style={rectangle, rounded corners=3pt, draw, font=\sffamily\scriptsize, align=center,
        minimum width=2.4cm, minimum height=0.8cm, inner sep=3pt},
    grp/.style={rectangle, rounded corners=5pt, draw=grisISO!60, dashed, inner sep=8pt},
    rel/.style={-{Stealth[length=2mm]}, semithick}]
    % Desarrollo
    \node[box, fill=c4Componente!45] (vite) at (0,1.4) {Vite\\\scriptsize :3000};
    \node[box, fill=c4Datos!22, draw=c4Datos!70!black] (spring1) at (0,0) {Spring\\\scriptsize :8080};
    \draw[rel] (vite) -- node[c4etiqueta, right]{\scriptsize proxy /api} (spring1);
    \begin{scope}[on background layer]
      \node[grp, fit=(vite)(spring1), label={[font=\sffamily\scriptsize\bfseries, text=colorPrimario]above:Desarrollo}] {};
    \end{scope}
    % Produccion
    \node[box, fill=c4Sistema, text=white] (jar) at (5.4,0.7) {JAR monolítico\\\scriptsize SPA + API :8080};
    \node[box, fill=c4Externo!30, draw=c4Externo!70!black, below=of jar] (tect) at (5.4,-0.5) {+ Tectonic};
    \draw[rel] (jar) -- (tect);
    \begin{scope}[on background layer]
      \node[grp, fit=(jar)(tect), label={[font=\sffamily\scriptsize\bfseries, text=colorPrimario]above:Producción (Docker)}] {};
    \end{scope}
\end{tikzpicture}
\caption{Topologías de los entornos de desarrollo y de producción.}
\label{fig:entornos}
\end{figure}

\subsection{Entorno de desarrollo: \textit{proxy} inverso de Vite}
\label{subsec:entorno_dev}

En desarrollo, el servidor de Vite (puerto \textbf{3000}) sirve la SPA con recarga en
caliente, mientras el backend Spring corre por separado en el puerto \textbf{8080}. Para
que el cliente llame a \pathbreak{/api} sin tropezar con CORS, Vite expone un
\textbf{\textit{proxy} inverso} que reenvía las rutas \pathbreak{/api}, \pathbreak{/auth} y
\pathbreak{/v1} al backend.

\begin{lstlisting}[language=JavaScript, caption={Proxy inverso de desarrollo (\texttt{vite.config.ts}).}, label={lst:vite-proxy}]
server: {
  port: 3000,
  proxy: {
    '/api':  { target: 'http://localhost:8080', changeOrigin: true },
    '/auth': { target: 'http://localhost:8080', changeOrigin: true },
    '/v1':   { target: 'http://localhost:8080', changeOrigin: true },
  },
}
\end{lstlisting}

\begin{NotaTecnica}
El \textit{proxy} establece además \texttt{Cross-Origin-Opener-Policy:
same-origin-allow-popups} para soportar el flujo de ventana emergente de OAuth
(capítulo~\ref{ch:frontend}, sección~\ref{subsec:oauth}, pág.~\pageref{subsec:oauth}) y una
lista de \texttt{allowedHosts} para los dominios de previsualización. Este \textit{proxy}
\textbf{solo existe en desarrollo}: en producción, el mismo origen sirve la SPA y la API,
por lo que CORS deja de ser un problema.
\end{NotaTecnica}

\subsection{Entorno de producción: contenedor Docker}
\label{subsec:entorno_prod}

En producción, el sistema se empaqueta como una \textbf{imagen Docker multi-etapa}
(capítulo~\ref{ch:backend}, sección~\ref{subsec:docker}, pág.~\pageref{subsec:docker}): una
etapa de compilación con el JDK~17 y una etapa de ejecución con el JRE~17 que incorpora el
compilador \textbf{Tectonic}. El contenedor expone el puerto \textbf{8080} y arranca con
\comando{java -jar app.jar}. La tabla~\ref{tab:prod-params} resume sus parámetros.

\vspace{0.2cm}
\noindent
\renewcommand{\arraystretch}{1.3}
\fontsize{8.5}{10.2}\selectfont\sffamily
\begin{tabularx}{\textwidth}{|L{3.6cm}|Y|}
\hline
\rowcolor{headerTabla}
\sffamily\textbf{Parámetro} & \sffamily\textbf{Valor} \\
\hline
Imagen base (ejecución) & \texttt{eclipse-temurin:17-jre-jammy}. \\ \hline
\rowcolor{grisTecnico}
Puerto expuesto & \texttt{8080}. \\ \hline
Perfil Spring activo & \texttt{prod} (vía \texttt{SPRING\_PROFILES\_ACTIVE}). \\ \hline
\rowcolor{grisTecnico}
Compilador LaTeX & Tectonic 0.15.0 (caché precalentada). \\ \hline
\textit{Entrypoint} & \texttt{java -jar app.jar}. \\ \hline
\end{tabularx}
\label{tab:prod-params}

\begin{AvisoNota}
La configuración por contenedor sigue el principio de \textbf{paridad
desarrollo-producción}: el mismo artefacto que se prueba localmente (con su compilador
LaTeX y su perfil de Spring) es el que se despliega, reduciendo la clase de errores que
solo aparecen ``en producción''. Los parámetros que varían entre entornos se inyectan
exclusivamente por variables de entorno, no por reconstrucción de la imagen, como detalla
la sección~\ref{sec:secretos} (pág.~\pageref{sec:secretos}).
\end{AvisoNota}
